\section{The quantized radiation field}\label{app:C.16}

The 2p state of hydrogen decays to 1s in \SI{1.60}{ns} by emitting one photon, and a silicon nucleus kicked by dark matter radiates in the same way, more weakly. \Cref{sec:12.6} computes such emission rates in a crystal. They need the electromagnetic field as a set of quantized oscillators, the coupling of a charge to it, the one-photon emission rate with its dipole limit, and the change of all three in a transparent medium of refractive index $n_r(\omega)$. We use Heaviside--Lorentz units with $e^2 = 4\pi\alpha$ (\cref{app:A.2}) and the box normalization of \cref{app:C.2}, with volume $\Vol$.

\subsection{The field as oscillators}

In the Coulomb gauge, $\nabla\cdot\vect{A} = 0$ and $A^0$ is the instantaneous Coulomb potential of the charges, already contained in the Coulomb interaction of \cref{sec:2.3}. The radiation field is the transverse vector potential. In the periodic box it is a sum over wavevectors $\vect{k}$ and two polarization vectors $\boldsymbol{\epsilon}_{\vect{k}j}$, $j = 1, 2$, orthonormal and perpendicular to $\vect{k}$, which makes $\nabla\cdot\vect{A} = 0$ term by term. Without sources each term obeys the wave equation with $\omega_k = k$:
\begin{equation*}
  \vect{A}(\vect{r}, t) = \sum_{\vect{k}j}\bigl[c_{\vect{k}j}\,\boldsymbol{\epsilon}_{\vect{k}j}\,\ee^{\ii\vect{k}\cdot\vect{r} - \ii\omega_kt} + \text{c.c.}\bigr] .
\end{equation*}
The fields are $\vect{E} = -\partial_t\vect{A}$ and $\vect{B} = \nabla\times\vect{A}$, so each mode contributes $\ii\omega_k c_{\vect{k}j}\boldsymbol{\epsilon}_{\vect{k}j}$ to $\vect{E}$ and $\ii\vect{k}\times\boldsymbol{\epsilon}_{\vect{k}j}c_{\vect{k}j}$ to $\vect{B}$, times the plane wave. The energy $\tfrac12\int\dd^3r\,(\vect{E}^2 + \vect{B}^2)$ is evaluated with $\int\dd^3r\,\ee^{\ii(\vect{k} - \vect{k}')\cdot\vect{r}} = \Vol\delta_{\vect{k}\vect{k}'}$. Terms that pair $\vect{k}$ with $-\vect{k}$ oscillate as $\ee^{\mp2\ii\omega_kt}$ and cancel between $\vect{E}^2$ and $\vect{B}^2$, because $(\vect{k}\times\boldsymbol{\epsilon})\cdot(\vect{k}\times\boldsymbol{\epsilon}') = k^2\,\boldsymbol{\epsilon}\cdot\boldsymbol{\epsilon}'$ for transverse vectors and $k^2 = \omega_k^2$. What remains is $\sum_{\vect{k}j}2\Vol\omega_k^2|c_{\vect{k}j}|^2$. A harmonic oscillator of frequency $\omega_k$ whose classical amplitude $\mathfrak{a}_{\vect{k}j}$ is normalized so that its energy is $\omega_k|\mathfrak{a}_{\vect{k}j}|^2$ matches this with $c_{\vect{k}j} = \mathfrak{a}_{\vect{k}j}/(2\Vol\omega_k)^{1/2}$. Quantization replaces $\mathfrak{a}_{\vect{k}j}$ by a boson annihilation operator $d_{\vect{k}j}$ with the commutators of \cref{eq:2.3-comm}:
\begin{equation}
  \vect{A}(\vect{r}) = \sum_{\vect{k}j}\frac{1}{(2\Vol\omega_k)^{1/2}}\bigl[d_{\vect{k}j}\,\boldsymbol{\epsilon}_{\vect{k}j}\,\ee^{\ii\vect{k}\cdot\vect{r}} + d^\dagger_{\vect{k}j}\,\boldsymbol{\epsilon}^*_{\vect{k}j}\,\ee^{-\ii\vect{k}\cdot\vect{r}}\bigr], \qquad
  H_{\mathrm{rad}} = \sum_{\vect{k}j}\omega_k\bigl(d^\dagger_{\vect{k}j}d_{\vect{k}j} + \tfrac12\bigr) .
  \label{eq:C.16-A}
\end{equation}
The eigenstates of $H_{\mathrm{rad}}$ are the photon number states $|\{n_{\vect{k}j}\}\rangle$, and by \cref{eq:2.3-boson} the operator $d^\dagger_{\vect{k}j}$ adds a photon with amplitude $(n_{\vect{k}j} + 1)^{1/2}$. Emission into a mode that already holds $n$ photons is therefore $n + 1$ times faster than into an empty one; for $\omega \gg T$ the thermal occupation $n_{\mathrm{B}}(\omega)$ is negligible and only spontaneous emission, the 1, remains. Complex polarization vectors describe circular polarization, and the formulas below hold for them.

\subsection{Minimal coupling}

A non-relativistic particle of mass $m$ and charge $e_{\mathrm{c}}$, with $e_{\mathrm{c}} = -e$ for an electron and $e_{\mathrm{c}} = Z_\kappa e$ for a bare nucleus, couples to $\vect{A}$ by \emph{minimal coupling}, $\vect{p} \to \vect{p} - e_{\mathrm{c}}\vect{A}$. For the electron this is $\vect{p} + e\vect{A}$, the non-relativistic limit of \cref{eq:C.15-HD}. Expanding the kinetic energy,
\begin{equation}
  \frac{(\vect{p} - e_{\mathrm{c}}\vect{A})^2}{2m} = \frac{p^2}{2m} - \frac{e_{\mathrm{c}}}{m}\,\vect{A}\cdot\vect{p} + \frac{e_{\mathrm{c}}^2}{2m}\,\vect{A}^2 ,
  \label{eq:C.16-minimal}
\end{equation}
where $\vect{p}\cdot\vect{A} = \vect{A}\cdot\vect{p} - \ii\nabla\cdot\vect{A} = \vect{A}\cdot\vect{p}$ in the Coulomb gauge. The term linear in $\vect{A}$ creates or absorbs one photon; the quadratic term acts at first order only in two-photon processes, such as scattering, and drops out of one-photon emission. For several particles the interaction is the sum over particles, and for a prescribed classical current $\vect{J}(\vect{r}, t)$ it is $H' = -\int\dd^3r\,\vect{J}\cdot\vect{A}$, the form taken by the linear term when the particle velocity $\vect{p}/m$ is a given function of time.

\subsection{The one-photon emission rate}

Let the matter make a transition $|i\rangle \to |f\rangle$ with energy release $\omega = E_i - E_f > 0$ and emit one photon into the empty field. Only the $d^\dagger$ part of \cref{eq:C.16-A} contributes, and the matrix element of the linear term of \cref{eq:C.16-minimal} is
\begin{equation*}
  \langle f; 1_{\vect{k}j}|H'|i; 0\rangle = -\frac{e_{\mathrm{c}}}{m\,(2\Vol\omega_k)^{1/2}}\,\boldsymbol{\epsilon}^*_{\vect{k}j}\cdot\vect{M}_{fi}, \qquad
  \vect{M}_{fi} = \langle f|\ee^{-\ii\vect{k}\cdot\vect{r}}\,\vect{p}|i\rangle .
\end{equation*}
The golden rule \cref{eq:2.1-golden}, summed over the photon modes with $\sum_{\vect{k}} \to \Vol\int\dd^3k/(2\pi)^3 = \Vol\int\omega_k^2\dd\omega_k\,\dd\Omega_{\hat{\vect{k}}}/(2\pi)^3$, gives the rate per solid angle of emission
\begin{equation}
  \frac{\dd\varGamma}{\dd\Omega_{\hat{\vect{k}}}} = \frac{e_{\mathrm{c}}^2\,\omega}{8\pi^2m^2}\sum_j\bigl|\boldsymbol{\epsilon}^*_{\vect{k}j}\cdot\vect{M}_{fi}\bigr|^2, \qquad |\vect{k}| = \omega .
  \label{eq:C.16-rate}
\end{equation}
The prefactor $\omega/(8\pi^2)$ collects $2\pi$ from the golden rule, $\omega^2/(2\pi)^3$ from the density of modes after the energy delta function is integrated, and $1/(2\omega)$ from the field normalization.

\emph{Dipole limit.} When the emitting system is small compared with the photon wavelength, $kr \ll 1$, the factor $\ee^{-\ii\vect{k}\cdot\vect{r}}$ is replaced by 1. A photon of \SI{10}{eV} has $k = \SI{5.1e-3}{\angstrom^{-1}}$, so across an atom $kr \approx 0.005$. If the unperturbed Hamiltonian is $H_0 = p^2/2m + U(\vect{r})$, then $[H_0, \vect{r}] = -\ii\vect{p}/m$, and $\langle f|\vect{p}|i\rangle = \ii m(E_f - E_i)\vect{r}_{fi} = -\ii m\omega\,\vect{r}_{fi}$, with $\vect{r}_{fi} = \langle f|\vect{r}|i\rangle$. The sum over the two polarizations is $\sum_j|\boldsymbol{\epsilon}^*_j\cdot\vect{d}|^2 = |\vect{d}|^2 - |\hat{\vect{k}}\cdot\vect{d}|^2$ for any vector $\vect{d}$, by completeness of $\{\boldsymbol{\epsilon}_1, \boldsymbol{\epsilon}_2, \hat{\vect{k}}\}$, and its integral over directions is $4\pi|\vect{d}|^2(1 - \tfrac13)$ (\cref{app:C.3}). The total rate is
\begin{equation}
  \varGamma_{i\to f} = \frac{e_{\mathrm{c}}^2\,\omega^3}{3\pi}\,|\vect{r}_{fi}|^2 = \frac43\,\alpha\,\omega^3\,|\vect{r}_{fi}|^2 \qquad (|e_{\mathrm{c}}| = e) .
  \label{eq:C.16-dipole}
\end{equation}
For hydrogen 2p $\to$ 1s, $\omega = \tfrac34\times\SI{13.606}{eV}$ and $|\vect{r}_{fi}|^2 = (2^{15}/3^{10})\,a_0^2$ for each $m$ substate; the lifetime $1/\varGamma$ is \SI{1.60}{ns}, the measured value. A transition with $\vect{r}_{fi} = 0$ proceeds through the next terms of the expansion of $\ee^{-\ii\vect{k}\cdot\vect{r}}$, magnetic-dipole and electric-quadrupole, smaller by $(kr)^2$.

\emph{Emission by a classical current.} For a prescribed current, first-order perturbation theory gives the amplitude of one photon in mode $\vect{k}j$ as $c_{\vect{k}j} = -\ii\int\dd t\,\langle 1_{\vect{k}j}|H'(t)|0\rangle = \ii(2\Vol\omega_k)^{-1/2}\,\boldsymbol{\epsilon}^*_{\vect{k}j}\cdot\tilde{\vect{J}}(\vect{k}, \omega_k)$, with $\tilde{\vect{J}}(\vect{k}, \omega) = \int\dd t\,\dd^3r\,\ee^{\ii\omega t - \ii\vect{k}\cdot\vect{r}}\,\vect{J}(\vect{r}, t)$ in the conventions of \cref{app:A.3}. Summing $|c_{\vect{k}j}|^2$ over modes as above, the number of photons emitted per unit frequency and solid angle is
\begin{equation}
  \frac{\dd N_\gamma}{\dd\omega\,\dd\Omega_{\hat{\vect{k}}}} = \frac{\omega}{16\pi^3}\sum_j\bigl|\boldsymbol{\epsilon}^*_{\vect{k}j}\cdot\tilde{\vect{J}}(\vect{k}, \omega)\bigr|^2 .
  \label{eq:C.16-classical}
\end{equation}
A charge at rest that is suddenly set moving with velocity $\vect{v}$ at $t = 0$ carries $\vect{J} = e_{\mathrm{c}}\vect{v}\,\delta^3(\vect{r} - \vect{v}t)\,\Theta(t)$, so $\tilde{\vect{J}} = e_{\mathrm{c}}\vect{v}\int_0^\infty\dd t\,\ee^{\ii(\omega - \vect{k}\cdot\vect{v})t} = \ii e_{\mathrm{c}}\vect{v}/(\omega - \vect{k}\cdot\vect{v} + \ii0^+) \simeq \ii e_{\mathrm{c}}\vect{v}/\omega$ for $v \ll 1$. With the polarization sum and angular integral of the dipole case,
\begin{equation}
  \frac{\dd N_\gamma}{\dd\omega} = \frac{e_{\mathrm{c}}^2\,v^2}{6\pi^2\,\omega} = \frac{2\alpha}{3\pi}\,\frac{v^2}{\omega} \qquad (|e_{\mathrm{c}}| = e) .
  \label{eq:C.16-soft}
\end{equation}
The energy radiated per unit frequency, $\omega\,\dd N_\gamma/\dd\omega$, is flat. The kick may be treated as sudden when it lasts less than $1/\omega$. For a nucleus bound in a crystal the restoring forces and the response of the electrons modify \cref{eq:C.16-soft} at the frequencies of the lattice and of the electrons, which \cref{sec:12.6} includes through the dynamical Born charge.

\subsection{A photon mode in a medium}

Inside a transparent medium with relative permittivity $\varepsilon_r(\omega) = n_r(\omega)^2$ and no magnetic response, the photon is a mode of the macroscopic field. The dispersion relation becomes $k = n_r(\omega)\,\omega$, and the number of modes per unit frequency grows: $\Vol\,k^2\,\dd k/(2\pi)^3 = \Vol\,n_r^2\,n_{\mathrm{grp}}\,\omega^2\dd\omega/(2\pi)^3$, with $n_{\mathrm{grp}} = \dd k/\dd\omega = n_r + \omega\,\dd n_r/\dd\omega$ the group index. The energy stored in a mode of given amplitude grows as well. The time-averaged energy of a monochromatic field in a transparent dispersive medium is $\tfrac14[\dd(\omega\varepsilon_r)/\dd\omega\,|\vect{E}_0|^2 + |\vect{B}_0|^2]$ per unit volume~\cite{jackson1999classical}. With $|\vect{E}_0| = \omega|\vect{A}_0|$, $|\vect{B}_0| = k|\vect{A}_0| = n_r\omega|\vect{A}_0|$ and $\dd(\omega\varepsilon_r)/\dd\omega = n_r^2 + 2n_r\omega\,\dd n_r/\dd\omega$, the bracket is $\omega^2|\vect{A}_0|^2\cdot2n_rn_{\mathrm{grp}}$, larger than in vacuum by the factor $n_rn_{\mathrm{grp}}$. Repeating the quantization, the normalization $(2\Vol\omega)^{-1/2}$ of \cref{eq:C.16-A} becomes $(2\Vol\omega\,n_rn_{\mathrm{grp}})^{-1/2}$; for a medium without dispersion, $n_{\mathrm{grp}} = n_r$, this is the quantization of Ref.~\cite{glauber1991quantum}.

In a rate, the squared normalization contributes $1/(n_rn_{\mathrm{grp}})$ and the density of modes $n_r^2n_{\mathrm{grp}}$. The group index cancels, and
\begin{equation}
  \varGamma^{\mathrm{medium}} = n_r(\omega)\,\varGamma^{\mathrm{vac}}, \qquad
  \frac{\dd N_\gamma^{\mathrm{medium}}}{\dd\omega} = n_r(\omega)\,\frac{\dd N_\gamma^{\mathrm{vac}}}{\dd\omega},
  \label{eq:C.16-medium}
\end{equation}
for the same matrix element or current. The medium must be transparent at $\omega$, so that $\varepsilon_r$ is real and the photon is a propagating mode; near an absorption line the photon is not a well-defined final state. And the emitter must couple to the macroscopic field: the local field at an atom differs from the macroscopic one by a local-field factor, which \cref{eq:C.16-medium} sets to 1. In the dipole factor $\ee^{-\ii\vect{k}\cdot\vect{r}}$ the wavevector becomes $n_r\omega$, which leaves the dipole limit intact for $n_r$ of a few.
